\subsection{The channels $(b,\gamma^{r})$ and $(\gamma^{r},b)$}%
\label{sec:superspace:bgamma}

We now compute the mixed vector--antichiral family: the six ordered channels $(b,\gamma^{r})$ and $(\gamma^{r},b)$,
$r=1,2,3$, of the master identity~\eqref{eq:ss:master}. The two ordered insertions are
\begin{equation}\label{eq:bg:insert}
I_{b\gamma^{r}}=\nabla_{-}\big(b^{A}\gamma^{rB}\big), \qquad I_{\gamma^{r}b}=\nabla_{-}\big(\gamma^{rA}b^{B}\big) .
\end{equation}
Both letters are Grassmann-even and $\nabla_{-}\gamma^{r}=0$, so in each insertion the outer derivative marks only
the $b$ line; every nonlinear occurrence of the covariantized source is nevertheless retained. The two directions
are computed independently below --- no graph coefficient is copied from one ordering to the other. As in the seed
channel, the graph weights are assembled with the unit-normalized field-strength insertion $X:=\nabla_{+}\cW_{+}$;
the letter normalization $b=-(i/\sqrt2)X$ is restored only in the final step.

\paragraph{1. Graph inventory.} The covariant completion of the insertion and its $\cD_{-}$-words are
\begin{equation}\label{eq:bg:source}
X_{c}=X_{1}+gX_{2}+g^{2}X_{3}, \qquad N_{0}=\cD_{-}X_{1}, \quad N_{1}=\cD_{-}X_{2}+[\Gamma_{1},X_{1}], \quad N_{2}=\cD_{-}X_{3}+[\Gamma_{1},X_{2}]+[\Gamma_{2},X_{1}],
\end{equation}
where $\nabla_{-}=\cD_{-}+g\Gamma_{1}+g^{2}\Gamma_{2}$ is the connection expansion of the semi-chiral derivative,
while $\gamma^{r}=\widetilde\Phi^{r}$ has no nonlinear completion. The connected order-$g^{2}$ orbit is
\begin{equation}\label{eq:bg:orbit}
\Gamma^{(1)}_{g^{2}} =\langle I_{2}\rangle_{0} -\frac1\h\langle I_{1}S_{3}\rangle_{0,c} -\frac1\h\langle I_{0}S_{4}\rangle_{0,c} +\frac1{2\h^{2}}\langle I_{0}S_{3}S_{3}\rangle_{0,c} ,
\end{equation}
with $I_{n}$ the source term at order $g^{n}$, $S_{3}$ and $S_{4}$ the cubic and quartic action vertices, and
$\langle\;\rangle_{0,c}$ the connected Gaussian average. The admitted parent triangle skeletons are the
gauge--matter class $T_{GM}^{(r)}=(I,G,M_{r})$ and the two-matter class $T_{MM}^{(r,r)}=(I,M_{r},M_{r})$, with
$G=\mathcal V_{VVV}$ the cubic gauge vertex and $M_{r}=\mathcal V_{\widetilde\Phi_{r}V\Phi_{r}}$ the matter--gauge
cubic; the superpotential class $T_{MH}=(I,M_{r},H_{\widetilde\Phi^{3}})$ contributes two complement-flavor routes
per direction, both exactly zero by
\eqref{eq:bg:tmh} below, and there is no all-gauge parent for this letter pair. Per ordered direction the route census is
\begin{equation}\label{eq:bg:census}
N_{TGM}=6, \qquad N_{TMM}=1, \qquad N_{TMH}=2 ,
\end{equation}
the six gauge--matter routes being the six permutations of the three internal edges $(r_{0},r_{1},r_{2})$, each
with two derivative placements times two Euclidean chiralities of the matter line. Figure~\ref{fig:bg:triangle}
shows the ordered triangle in the two-matter topology.

\begin{figure}[t]
\centering
\begin{tikzpicture}[scale=1.15]
  \coordinate (I)  at (90:2.0);
  \coordinate (ML) at (210:2.0);
  \coordinate (MR) at (330:2.0);
  \draw[sV]    (I)  -- (ML) node[midway,left=3pt]  {$r_{0}$};
  \draw[sPhi]  (ML) -- (MR) node[midway,below=3pt] {$r_{1}$};
  \draw[sPhit] (MR) -- (I)  node[midway,right=3pt] {$r_{2}$};
  \draw[sPhit] (ML) -- ++(180:1.5) node[left]  {$\widetilde\Phi^{r}\ (p)$};
  \draw[sleg] (MR) -- ++(0:1.5)   node[right] {$\partial_{\dot\alpha}c\ (q)$};
  \node[sinsert,label={90:$\nabla_{-}\bigl(b^{A}\gamma^{rB}\bigr)$}] at (I) {};
  \node[svertex,label=below left:{$M_{r}$}]  at (ML) {};
  \node[svertex,label=below right:{$M_{r}$}] at (MR) {};
\end{tikzpicture}
\caption{The ordered triangle of the $(b,\gamma^{r})$ family in the two-matter topology $T_{MM}^{(r,r)}$. The filled
insertion node carries the letter operator $\nabla_{-}(b^{A}\gamma^{rB})$ with $b=-(i/\sqrt2)\nabla_{+}\cW_{+}$;
the internal edges are the source-$b$ vector line $r_{0}$ (wavy), the chiral matter bridge $r_{1}$, and the
source-$\gamma$ reversed-chiral line $r_{2}$, with $r_{1}-r_{0}=p$. The external antichiral letter
$\gamma^{r}=\widetilde\Phi^{r}$ (momentum $p$) enters at the left matter vertex, the vertex where the vector
edge $r_{0}$ attaches (each matter cubic vertex has exactly one vector slot); the vector output letter
$\partial_{\dot\alpha}c$ (momentum $q$) enters at the right matter vertex. The gauge--matter topology
$T_{GM}^{(r)}$ replaces the left matter vertex by the cubic gauge vertex $G$ (the edges incident to $G$ are
vector lines); its six ordered routings of $(r_{0},r_{1},r_{2})$ are the six permutation routes of the text, and
the reversed channel $(\gamma^{r},b)$ has the same figure with the source attachments exchanged.}
\label{fig:bg:triangle}
\end{figure}

\paragraph{2. Wick contractions.} Every labeled Wick contraction carries weight $1$: the $1/2!$ of the two cubic action
copies is canceled by the two labeled action orderings,
\begin{equation}\label{eq:bg:wick}
w_{\mathrm{Wick}}=\frac{1}{2!}\,(1+1)=1 ,
\end{equation}
and no residual symmetry factor occurs anywhere in the family. In the triangle term $\langle
I_{0}S_{3}S_{3}\rangle_{0,c}$ the two external ports attach to two distinct cubic vertices and the three internal
edges are paired among the vertex slots, giving the six $T_{GM}$ routings, the single $T_{MM}$ routing (two matter
propagators, hence two chiral projectors), and the two $T_{MH}$ routings. In the nonlinear-source bubble $\langle
I_{1}S_{3}\rangle_{0,c}$ the two internal legs of the connection word $[\Gamma_{1},X_{1}]$ pair with each other
through the loop, and the cubic matter vertex ties the loop to the external $\gamma^{r}$ leg. The quartic seagull
$\langle I_{0}S_{4}\rangle_{0,c}$ is a single graph: differentiating
\begin{equation}\label{eq:bg:seagull}
[\cV,[\cV,\Phi]]=\cV^{2}\Phi-2\cV\Phi\cV+\Phi\cV^{2}
\end{equation}
with respect to one internal and one external $\cV$ places both orderings inside one seagull vertex; they are not
two Feynman graphs. The remaining source tadpole $\langle I_{2}\rangle_{0}$ is a single coincident vector loop. The
full color contraction of the family --- two matter vertices, two matter inverse metrics, and the complete vector
covariance --- reduces, with $[T_{A},T_{B}]=if_{AB}{}^{C}T_{C}$ and $\operatorname{tr}(T_{A}T_{B})=\delta_{AB}/2$,
to
\begin{equation}\label{eq:bg:color}
C_{\mathrm{full}} \;\propto\; \sum_{X}\operatorname{tr}\big(T_{D}[T_{A},T_{X}]\big) \operatorname{tr}\big(T_{X}[T_{E},T_{B}]\big) \;\longrightarrow\; +\,f^{ACD}f^{BCE} ,
\end{equation}
the numerical proportionality being carried into the assembled prefactors of paragraph 3. This is the adjudicated
signed color word of the family: the same tensor unit as in the master identity~\eqref{eq:ss:master}.

\paragraph{3. Amplitude assembly.} The propagators entering the orbit are the vector line and the two matter
orientations,
\begin{equation}\label{eq:bg:props}
-\frac{\h}{\BoxE}, \qquad +\frac{\h\,\bcD^{2}\cD^{2}}{16\,\BoxE}, \qquad +\frac{\h\,\cD^{2}\bcD^{2}}{16\,\BoxE} ;
\end{equation}
each matter vertex contributes $+\sqrt2\,g/\h$ with measure $1/4$ per vertex, and the canonical source
normalization is $-1/(4\sqrt2)$. The two-matter orbit therefore assembles to the exact pre-$D$ factor
\begin{equation}\label{eq:bg:tmmpre}
\left(-\frac{1}{4\sqrt2}\right)\left(\frac{\sqrt2\,g}{\h}\right)^{2} \left[-\h\left(\frac{\h}{16}\right)^{2}\right]=\frac{\sqrt2\,\h g^{2}}{1024} ,
\end{equation}
while each gauge--matter route carries, per Euclidean chirality,
\begin{equation}\label{eq:bg:tmpre}
\mathrm{preD}_{-}=+\frac{\h g^{2}}{8192\sqrt2}, \qquad \mathrm{preD}_{+}=-\frac{\h g^{2}}{8192\sqrt2} .
\end{equation}
After loop saturation the external legs carry residual operator words: the $\gamma^{r}$ leg is a plain antichiral
field (factor $1$), and the vector output leg collapses as
\begin{equation}\label{eq:bg:residual}
\cD^{2}\bcD_{\dot\alpha}\,u=-4\sqrt2\,D_{\dot\alpha}, \qquad \cD^{2}\bcD_{\dot\alpha}\,(\theta+\bar\theta\,\theta_{\dot\alpha})=-2 ,
\end{equation}
where $u$ is the formal leg slot and $D_{\dot\alpha}$ the unit vector output slot (the antichiral field strength
$\bcW_{\dot\alpha}$; the letter $\partial_{\dot\alpha}c=i\bcW_{\dot\alpha}$ is restored in paragraph 7). The
ordered output kernels are $(\partial_{\dot\alpha}Y)^{D}D^{E\dot\alpha}$ for the two-matter orbit and
$D_{\dot\alpha}^{D}(\partial^{\dot\alpha}Y)^{E}$ for the gauge--matter orbit, with $Y=\widetilde\Phi^{r}$.

\paragraph{4. The $D$-algebra chain.} The master $D$-word behind every route is
\begin{equation}\label{eq:bg:masterD}
\cD_{-}\cD_{+}\bcD^{2}\cD_{+} =8\,\bar\Box\,\cD_{+}-\frac12\,\cD_{+}\bcD^{2}\cD^{2} ,
\end{equation}
with $\bar\Box$ the d'Alembertian of the barred spin metric of \eqref{eq:ss:sigmaalgebra}; loop saturation uses
$\delta^{4}(\theta_{12})\cD^{2}\bcD^{2}\delta^{4}(\theta_{12})=16\,\delta^{4}(\theta_{12})$ as in the seed channel.
For the two-matter route the four-dimensional numerator before the metric split is
\begin{equation}\label{eq:bg:tmmnum}
\mathcal F^{(4)}_{b>\gamma} =8(\det r_{0})\,\mathcal R_{0}-8(\det r_{1})\,\mathcal J_{1}, \qquad \mathcal R_{0}=-128\,\mathfrak b(r_{1}), \qquad \mathcal J_{1}=-128\,\mathfrak b(r_{0}) ,
\end{equation}
where $r_{0}$ is the source-$b$ vector edge, $r_{1}$ the chiral bridge, $\det r_{e}$ the $2\times2$ determinant of
the edge momentum (kept symbolic; its index-level definition is not needed below), and $\mathfrak b(\cdot)$ the
external bilinear carrier. The collapsed neighbor of the same cut orbit, its migrated and transverse residuals, and
the projector-collapse check are
\begin{equation}\label{eq:bg:tmmcontact}
K_{\mathrm{neighbor}}=(64\,r_{0\dot2},\,-64\,r_{0\dot1}), \quad (-128\,r_{0\dot2},\,+128\,r_{0\dot1}), \quad \big(+128(p_{\dot2}-r_{0\dot2}),\,-128(p_{\dot1}-r_{0\dot1})\big), \quad 16\det(r_{1})\,K_{\mathrm{neighbor}}=-8\det(r_{1})\,\mathcal J_{1} ,
\end{equation}
with the longitudinal tag migrating the source vector $r_{0}$ into the adjacent chiral projector $r_{1}$; for the
reversed direction the inverse-edge rows read $+8\det(r_{2})\mathcal R_{1}$ (primary) and $-8\det(r_{1})\mathcal
R_{2}$ (neighbor after longitudinal integration by parts). The finite-Grassmann replay of the nonlinear contact
words returns the same polynomials in the reversed ports for the two directions; no reflected coefficient is
assumed. Route by route, the full-$d$ Schwinger contact identity
\begin{equation}\label{eq:bg:sdcontact}
\mathcal F^{(4)}_{\mathrm{route}} +\mathcal C^{(d)}_{\mathrm{source}} +\mathcal C^{(d)}_{\mathrm{neighbor}}=0 \qquad(\mul^{2}=0)
\end{equation}
holds before the evanescent remainder is extracted. The six gauge--matter routes give the following signed traces
(full, longitudinal, transverse; per Euclidean chirality $\pm$), reproduced identically at two independent
external-momentum samples:
\begin{center}
\begin{tabular}{ccrrrrrr}
\hline
route & permutation & $-$ full & $-$ long. & $-$ trans. & $+$ full & $+$ long. & $+$ trans.\\
\hline
001 & $(0,1,2)$ & $-512$  & $1024$  & $-1536$ & $1024$ & $1024$  & $0$\\
002 & $(0,2,1)$ & $-1024$ & $-1024$ & $0$     & $0$    & $-512$  & $512$\\
003 & $(1,0,2)$ & $-1024$ & $-1024$ & $0$     & $512$  & $-512$  & $1024$\\
004 & $(2,0,1)$ & $-512$  & $512$   & $-1024$ & $1024$ & $1024$  & $0$\\
005 & $(1,2,0)$ & $-1024$ & $-512$  & $-512$  & $0$    & $0$     & $0$\\
006 & $(2,1,0)$ & $0$     & $0$     & $0$     & $1536$ & $0$     & $1536$\\
\hline
\end{tabular}
\end{center}
The per-route coefficients of the master unit, $-3/8,-1/4,-3/8,-3/8,-1/4,-3/8$ in the order of the table, sum to
$-2$; combined with the signed color word $-f^{ACD}f^{BCE}$ carried by the gauge--matter orbit before assembly,
this supplies the $+2$ of paragraph 6.

\paragraph{5. Evanescent extraction.} On every marked edge $e$ of every route the two metric splits of
Section~\ref{sec:superspace:cutting} give
\begin{equation}\label{eq:bg:edge}
\frac{r_{e,d}^{2}\,\mathcal R_{e}}{D_{0}D_{1}D_{2}} -\frac{\mathcal R_{e}}{\prod_{j\neq e}D_{j}}=0, \qquad \frac{\bar r_{e}^{\,2}\,\mathcal R_{e}}{D_{0}D_{1}D_{2}} -\frac{\mathcal R_{e}}{\prod_{j\neq e}D_{j}} =\frac{\mul^{2}\,\mathcal R_{e}}{D_{0}D_{1}D_{2}} ,
\end{equation}
the numerator-weighted form of the exact cutting-failure identity~\eqref{eq:ss:cuttingfailure}; no graph-specific
$(4-d)$ factor is inserted. With the ordered Feynman parametrization of \eqref{eq:ss:feynman}, the numerator
integral is
\begin{equation}\label{eq:bg:i2}
I_{\ell^{2}}(\Delta) :=\mu^{2\epsilon}\!\int\!\frac{d^{d}\ell}{(2\pi)^{d}}\, \frac{\ell_{d}^{2}}{(\ell_{d}^{2}+\Delta)^{3}} =J_{2}-\Delta J_{3}
\end{equation}
in the notation of \eqref{eq:ss:Jn}, and the evanescent master integral of
\eqref{eq:ss:masterintegral}--\eqref{eq:ss:masterproof} reads
\begin{equation}\label{eq:bg:jmu2}
J_{\mu^{2}} :=\mu^{2\epsilon}\!\int\!\frac{d^{d}\ell}{(2\pi)^{d}}\, \frac{\mul^{2}}{(\ell_{d}^{2}+\Delta)^{3}} =\frac{4-d}{d}\,I_{\ell^{2}} =\frac{1}{32\pi^{2}} .
\end{equation}
For the two-matter route the contact sum rule and evanescent probe are
\begin{equation}\label{eq:bg:sumrule}
\mathcal F^{\mathrm{DRED}}_{b>\gamma} +\mathcal C^{(d)}_{0}+\mathcal C^{(d)}_{1} =8\mul^{2}\big(\mathcal R_{0}-\mathcal J_{1}\big) =-1024\,\mul^{2}\,\mathfrak b(p) ,
\end{equation}
collapsing to $(-1024\,p_{\dot2},\,+1024\,p_{\dot1})$ on the two dotted spinor components (the ordering convention of
the recorded orbit replay), with $r_{1}-r_{0}=p$ by momentum conservation, so that $\mathfrak b(r_{1})-\mathfrak
b(r_{0})=\mathfrak b(p)$. The reversed direction gives $8\mul^{2}(\mathcal R_{2}-\mathcal R_{1})$ with probe
$(-1024\,q_{\dot2},\,+1024\,q_{\dot1})$, and the zero-source-momentum check at $q=-p$ returns the same magnitude.
The two superpotential routes are exact zeros,
\begin{equation}\label{eq:bg:tmh}
T_{MH}^{(1)}=T_{MH}^{(2)}=0 ,
\end{equation}
each with vanishing longitudinal, transverse, and full projected polynomials.

\paragraph{6. Pole cancellation.} In the two-matter orbit the parent triangle and the seagull contribute, in their
$I_{\ell^{2}}$ parts,
\begin{equation}\label{eq:bg:tmmpole}
\Gamma^{\triangle}_{TMM}\Big|_{I_{\ell^{2}}}=-\frac8d\,I_{\ell^{2}}, \qquad \Gamma^{I_{0}S_{m4}}_{TMM}\Big|_{I_{\ell^{2}}}=+2\,I_{\ell^{2}} \quad\Longrightarrow\quad \Big(-\frac84+2\Big)I_{\ell^{2}}=0 \quad(d=4) .
\end{equation}
In DRED the same pair leaves
\begin{equation}\label{eq:bg:tmmdred}
\Gamma^{\mathrm{DRED}}_{TMM}\Big|_{\mathrm{local}} =\Big(2-\frac8d\Big)I_{\ell^{2}} =\frac{2d-8}{d}\,I_{\ell^{2}} =-2\,\frac{4-d}{d}\,I_{\ell^{2}} =-2J_{\mu^{2}} \quad\Longrightarrow\quad \h g^{2}\,(-2)J_{\mu^{2}} =-\frac{\h g^{2}}{16\pi^{2}} .
\end{equation}
The gauge--matter orbit pairs its parent triangle with both the nonlinear-source bubble and the seagull:
\begin{equation}\label{eq:bg:tmpole}
\Gamma^{\triangle}_{TGM}\Big|_{I_{\ell^{2}}}=+\frac8d\,I_{\ell^{2}}, \quad \Gamma^{I_{1}S_{m3}}_{TGM}\Big|_{I_{\ell^{2}}}=-I_{\ell^{2}}, \quad \Gamma^{I_{0}S_{m4}}_{TGM}\Big|_{I_{\ell^{2}}}=-I_{\ell^{2}} \quad\Longrightarrow\quad \Big(\frac84-1-1\Big)I_{\ell^{2}}=0 \quad(d=4) ,
\end{equation}
\begin{equation}\label{eq:bg:tmdred}
\Gamma^{\mathrm{DRED}}_{TGM}\Big|_{\mathrm{local}} =\Big(\frac8d-2\Big)I_{\ell^{2}} =\frac{8-2d}{d}\,I_{\ell^{2}} =+2\,\frac{4-d}{d}\,I_{\ell^{2}} =+2J_{\mu^{2}} \quad\Longrightarrow\quad \h g^{2}\,(+2)J_{\mu^{2}} =+\frac{\h g^{2}}{16\pi^{2}} .
\end{equation}
The remaining order-$g^{2}$ rows vanish: the $X_{3}$ source tadpole is scaleless,
\begin{equation}\label{eq:bg:tadpole}
\langle I_{2}\rangle_{0} \;\supset\; \int\!\frac{d^{d}k}{(2\pi)^{d}}\,\frac{(k^{2})^{n}}{k^{2}}=0 ;
\end{equation}
the superpotential routes vanish by \eqref{eq:bg:tmh}; and the gauge, gauge-fixing, Faddeev--Popov, and
Nielsen--Kallosh rows have zero color trace in this sector and contain no $J_{\mu^{2}}$ remainder. The
Euler-current and explicit-current rows of the outer reduction of the unit insertion
$X$, $\nabla_{-}X=-\nabla_{+}(\delta
S/\delta\cV)-2i\sqrt2\,(\beta_{s}\!\times\!\gamma^{s})$ (flavor-summed), remain separate until their ports, denominators,
and edge tags are matched:
\begin{equation}\label{eq:bg:current}
E_{\mathrm{cur}}:\ +2i\ \text{on}\ D_{0}D_{1}, \qquad X_{\mathrm{cur}}:\ -2i\ \text{on}\ D_{0}D_{1}, \qquad (+2i-2i)\,K_{\mathrm{cur}}(k)=0 ,
\end{equation}
(the $\pm2i$ are quoted in the recorded machine units of the current rows,
i.e. before the $\sqrt2$ of the letter normalization in the reduction above),
with no marked inverse edge on either row (selected square $\varnothing$); this pairing vanishes before loop
integration and does not produce the factor $1/2$ below.

\begin{remark}[The occurrence decomposition]\label{rem:bg:occurrence}
An earlier replay of this family transported both inverse-square rows $\bar r_{0}^{2}-r_{0,d}^{2}=\mul^{2}$ and
$\bar r_{1}^{2}-r_{1,d}^{2}=\mul^{2}$ to independent nonlinear descendants and recorded
$|\Gamma_{\mathrm{old}}|=4J_{\mu^{2}}$, twice the accepted value. The complete orbit shows that these are two
decompositions of one metric pole orbit, so the complete value is $|\Gamma_{\mathrm{complete}}|=2J_{\mu^{2}}$: the
factor $1/2$ is the occurrence-decomposition correction, not the current-pair cancellation of
\eqref{eq:bg:current}. An external review of this family proposing a different repair was rejected for an
arithmetic gap; the occurrence decomposition above is the adjudicated mechanism.
\end{remark}

\paragraph{7. Final result.} In the ordered basis
\begin{equation}\label{eq:bg:basis}
\big(e_{1},e_{2}\big)^{DE} :=\Big((\partial_{\dot\alpha}\gamma^{r})^{D}(\partial^{\dot\alpha}c)^{E},\; (\partial_{\dot\alpha}c)^{D}(\partial^{\dot\alpha}\gamma^{r})^{E}\Big) =\Big(\langle\gamma^{r},\partial c\rangle,\, \langle\partial c,\gamma^{r}\rangle\Big)^{DE}
\end{equation}
the two surviving orbits of paragraph 6 sum to the independent coefficient pairs
\begin{equation}\label{eq:bg:coeffs}
\mathbf c_{b>\gamma^{r}}=(-1,+1), \qquad \mathbf c_{\gamma^{r}>b}=(-1,+1)
\end{equation}
in units of $(\h g^{2}/16\pi^{2})\,f^{ACD}f^{BCE}$: for the forward direction the two-matter orbit contributes
$-1$ on $e_{1}$ and the gauge--matter orbit $+1$ on $e_{2}$; the reversed direction exchanges the orbit--word
assignment and returns the same pair. Equivalently, in the unit normalization $X=\nabla_{+}\cW_{+}$,
$Y=\widetilde\Phi^{r}$, $D_{\dot\alpha}=\bcW_{\dot\alpha}$ of paragraphs 1--6,
\begin{equation}\label{eq:bg:unitresult}
\nabla_{-}\big(X^{A}Y^{B}\big)\Big|_{\mathrm{1\text{-}loop}} =\frac{\h g^{2}}{16\pi^{2}}\,f^{ACD}f^{BCE}\, \Big[D_{\dot\alpha}^{D}\big(\partial^{\dot\alpha}Y\big)^{E} -\big(\partial_{\dot\alpha}Y\big)^{D}D^{E\dot\alpha}\Big] .
\end{equation}
Restoring the letters of Definition~\ref{def:ss:letters} --- one insertion factor $-i/\sqrt2$ from
$b=-(i/\sqrt2)X$, none from $\gamma^{r}=Y$, and $D_{\dot\alpha}=-i\,\partial_{\dot\alpha}c$ from
$\partial_{\dot\alpha}c=i\bcW_{\dot\alpha}$ on each output word ---
\begin{equation}\label{eq:bg:conversion}
\nabla_{-}\big(b^{A}\gamma^{rB}\big) =-\frac{i}{\sqrt2}\cdot\frac{\h g^{2}}{16\pi^{2}}\,f^{ACD}f^{BCE}\, (-i)\,\big\{\langle\partial c,\gamma^{r}\rangle -\langle\gamma^{r},\partial c\rangle\big\}^{DE} ,
\end{equation}
and $(-i/\sqrt2)(-i)=-1/\sqrt2$ gives the channel formula
\begin{equation}\label{eq:bg:final}
\boxed{\; \nabla_{-}\big(b^{A}\gamma^{rB}\big) =\nabla_{-}\big(\gamma^{rA}b^{B}\big) =-\frac{\h g^{2}}{16\sqrt2\,\pi^{2}}\,f^{ACD}f^{BCE}\, \big\{\langle\partial c,\gamma^{r}\rangle -\langle\gamma^{r},\partial c\rangle\big\}^{DE} .\;}
\end{equation}
This is row \eqref{eq:ss:ch-bgamma} of the channel table. The two orderings are equal for a one-line reason: both
letters are Grassmann-even, so the graded Leibniz rule gives no Koszul sign under exchange, and
$\nabla_{-}\gamma^{r}=0$, so both insertions expand into the same source orbit with reversed attachments --- and
the two independent replays return the same ordered bracket, as \eqref{eq:bg:coeffs} records.

\paragraph{8. Holomorphic-twist comparison.} Budzik and Kulp compute the one-loop shift of the holomorphic supercharge
on the same letter pair (their equation (3.67))~\cite{BK}:
\begin{equation}\label{eq:bg:bk}
Q_{1}\big(b^{A}(\gamma^{I})^{B}\big) =\kappa^{2}f_{ACD}f_{BCE}\left[ \partial_{\dot\alpha}c^{D}\,\partial^{\dot\alpha}(\gamma^{I})^{E} -\partial_{\dot\alpha}(\gamma^{I})^{D}\,\partial^{\dot\alpha}c^{E}\right] .
\end{equation}
The output word is exactly $\langle\partial c,\gamma^{r}\rangle-\langle\gamma^{r},\partial c\rangle$: translating
\eqref{eq:bg:bk} into the present alphabet, the coefficient pair in the ordered basis \eqref{eq:bg:basis} is $(-1,+1)$,
in word-by-word agreement with \eqref{eq:bg:coeffs}, and the raw ratio to the printed zero-shift row is $1$. On the
holomorphic-twist side the row is obtained by the $\mathfrak{psl}(3|3)$ descent
$Q_{1}(b\gamma^{I})=-Q_{+}^{I}Q_{1}(bc)$ from the $(b,c)$ channel.

\begin{remark}[Status of the comparison]\label{rem:bg:ht}
The reversed row $Q_{1}((\gamma^{I})^{A}b^{B})$ is not printed in~\cite{BK}; the $(\gamma^{r},b)$ channel above is
an independent computation, not a quotient of the forward row. The comparison of this paragraph is an after-check
only: no coefficient of~\cite{BK} enters the derivation in paragraphs 1--7, and the agreement is a cross-check, not
an internal completeness proof of the graph census.
\end{remark}

\begin{remark}[Recorded granularity]\label{rem:bg:gaps}
The sources record the stage outputs of the $D$-algebra --- traces, residuals, probes, and the table above --- but
not the index-level intermediate spinor manipulations of the seven surviving routes; accordingly
\eqref{eq:bg:tmmnum}--\eqref{eq:bg:sdcontact} display the recorded outputs of each step rather than unattributed
intermediate lines, and Figure~\ref{fig:bg:triangle} is drawn from the recorded edge routings and the vertex
slot structure.
\end{remark}

